\begin{lemma}[Transporting the complete minimum context]
Fix a hierarchy counterexample $C$ on the exact tree $T$ with its minimum
context. Let $\mathcal C$ be a corresponding certificate under $\sigma$
on the same exact $T$, and let $C'$ be a hierarchy counterexample on a
base of its replacement ETT. Suppose $C.\mathit{good}$ and
$C'.\mathit{good}$ satisfy SameGoodHierarchy. Then $C'$ admits a minimum
context relative to the replacement input and generating coloring
$\sigma$. In particular all six lower-rung assertions, both minimality
conditions, and the $2n+11$ terminal missing-color bound are retained.
\end{lemma}
\begin{proof}
We supply all six fields of
\noderef{HierarchyMinimumCounterexampleContext}; its lower-rung field
itself contains all six technical assertions. Write $n$ for the original
rung and $n'$ for the replacement rung.
By \noderef{CorrespondingExtendedTashkinovTree}, $n'=n$, the replacement
ETT's series is the displayed replacement series, and $\sigma$ is the
coloring at its terminal index. The properness and one-uncoloredness
fields of that series at this index thus give the first two context
fields, using \noderef{IsProperEdgeColoring} and
\noderef{IsOneEdgeUncolored}.

For the lower-rung field, fix any certified positive-rung lower ETT, with
arbitrary input and generating coloring, and assume its rung $s$ is
strictly less than $n'$ and it has maximum property. Rung equality makes
$s<n$. Apply the original context's lower-rung field to this actual lower
ETT and its three stated premises. Its conclusion is the entire six-part
conjunction: elementarity; final-closure interchangeability under stable
colorings; the two PE-support stability and disjointness clauses; the
terminal PE component assertion; the defective-boundary vertex order
assertion; and terminal-stable mod-coloring. We return that conjunction
unchanged, not merely its first clause. This verifies the complete
lower-rung field for the new context.

The divider counts are equal by \noderef{SameGoodHierarchy}. For the
fewer-divider field, fix any maximum ETT of rung $n'$ and any good hierarchy
on it having fewer dividers than $C'$. The rung and divider equalities make
it a maximum ETT of rung $n$ with fewer dividers than $C$. The old
fewer-divider field supplies its elementarity under its own generating
coloring; return that assertion.

Apply \noderef{P2752_T2_HierarchyMeasures} to the same-stage clause of
the given same-good relation. Its starred-tree equality, path-number
equality and last ordinary-stage equality, substituted in
\noderef{HierarchyCounterexampleMeasure}, give exact equality of both
coordinates for $C,C'$. Now fix any counterexample of rung $n'$ and
the same divider count as $C'$. It has the original rung and divider count
as well. The old lexicographic-minimality field says its measure is not
less than that of $C$. Replacing that measure by the equal measure of
$C'$ gives exactly the new lexicographic-minimality field, with the order
of \noderef{HierarchyCounterexampleLess} unchanged.

For the large missing-set field, $n\ge1$ by
\noderef{ExtendedTashkinovTree}. The positive slot whose successor index
is $n$ therefore exists, even when $n=1$. At this slot, the ladder-tree
agreement of the corresponding certificate identifies the new terminal
ladder with the old $T_n$. Its replacement-stability field gives equality
of missing sets at every vertex of this old $T_n$, by the vertex clause
of \noderef{StableColoring}. The old generating-color equation and the
new terminal-color equation rewrite the two colorings in this equality
as $\varphi$ and $\sigma$. Taking the union over the exact vertex set,
as in \noderef{TreeMissingColors}, gives
\[
 \overline\sigma(T'_n)=\overline\varphi(T_n).
\]
Substitute this set equality and $n'=n$ into the original context's
$2n+11$ lower bound. This proves precisely the new large missing-set field.
All context fields have now been constructed for the supplied $C'$.
\end{proof}
